% ============================================================================
% Documentation technique de la pipeline d'evaluation du potentiel
% photovoltaique des toitures (nouvelle pipeline).
% Compiler avec pdflatex (2 passes pour la table des matieres).
% ============================================================================
\documentclass[11pt,a4paper]{article}

\usepackage[utf8]{inputenc}
\usepackage[T1]{fontenc}
\usepackage{lmodern}
\usepackage[french]{babel}
\usepackage[margin=2.4cm]{geometry}
\usepackage{amsmath,amssymb}
\usepackage{booktabs}
\usepackage{xcolor}
\usepackage{listings}
\usepackage{enumitem}
\usepackage[colorlinks=true, linkcolor=blue!50!black, urlcolor=blue!50!black,
            pdftitle={Pipeline PV toitures -- documentation technique}]{hyperref}

% ----------------------------------------------------------------------------
% Mise en forme des extraits de code Python
% ----------------------------------------------------------------------------
\definecolor{codegray}{rgb}{0.97,0.97,0.97}
\definecolor{codegreen}{rgb}{0.0,0.45,0.0}
\definecolor{codeblue}{rgb}{0.1,0.1,0.6}
\lstset{
  language=Python,
  basicstyle=\ttfamily\small,
  keywordstyle=\color{codeblue}\bfseries,
  commentstyle=\color{codegreen}\itshape,
  stringstyle=\color{purple},
  backgroundcolor=\color{codegray},
  showstringspaces=false,
  breaklines=true,
  frame=single,
  framerule=0.3pt,
  numbers=none,
  extendedchars=true,
  literate={é}{{\'e}}1 {è}{{\`e}}1 {ê}{{\^e}}1 {ë}{{\"e}}1
           {à}{{\`a}}1 {â}{{\^a}}1 {î}{{\^i}}1 {ï}{{\"i}}1
           {ô}{{\^o}}1 {ù}{{\`u}}1 {û}{{\^u}}1 {ç}{{\c c}}1
           {É}{{\'E}}1 {È}{{\`E}}1 {À}{{\`A}}1 {°}{{$^\circ$}}1,
}

% Raccourci pour les noms de fichiers / fonctions dans le texte
\newcommand{\fichier}[1]{\texttt{#1}}
\newcommand{\fonction}[1]{\texttt{#1}}

% filet de securite : evite les rares debordements de lignes (mots en tt longs)
\setlength{\emergencystretch}{3em}

\title{Évaluation du potentiel photovoltaïque des toitures\\
       \large Documentation technique complète de la pipeline}
\author{Stage M1 -- Université de Rennes}
\date{Juin 2026}

\begin{document}

\maketitle

\begin{abstract}
Ce document décrit en détail la pipeline de calcul du potentiel photovoltaïque
(PV) des toitures développée pendant le stage. Il présente brièvement le but du
projet et l'architecture générale, puis détaille \emph{chaque fonction} du code~:
ce qu'elle fait, comment elle est construite, pourquoi ces choix, et les concepts
(géomatiques, solaires, numériques) nécessaires pour la comprendre. Il se termine
par la liste de ce qu'il reste à développer. La méthode suit l'étude suisse de
Buffat \emph{et al.} (2018), \emph{A scalable method for estimating rooftop solar
irradiation potential over large regions}, Applied Energy 216.
Le document est écrit pour un lecteur qui ne connaît ni les bibliothèques
utilisées (rasterio, geopandas, pvlib\ldots) ni les données (LiDAR HD, BD TOPO,
PVGIS), mais qui veut pouvoir lire et modifier le code.
\end{abstract}

\tableofcontents
\newpage

% ============================================================================
\section{Le projet en un coup d'œil}
% ============================================================================

\subsection{Objectif}

Le but est de construire un outil capable d'estimer, \textbf{toiture par
toiture}, la quantité d'énergie solaire annuelle (et mensuelle) qu'un toit
reçoit, puis à terme l'électricité qu'une installation photovoltaïque y
produirait. L'outil doit fonctionner \textbf{à grande échelle}~: un bâtiment,
une commune, un département, et idéalement la France entière. Cette contrainte
d'échelle conditionne toute l'architecture~: on ne peut pas se permettre des
calculs lourds par bâtiment, il faut pré-calculer tout ce qui peut l'être et ne
faire au niveau du pixel de toiture que des opérations très simples.

\subsection{La méthode de référence : Buffat et al. (2018)}

L'étude suisse de Buffat \emph{et al.} a montré qu'on peut estimer le potentiel
solaire de plus de 3~millions de bâtiments (toute la Suisse, à 0,5~m de
résolution) avec une bonne précision (validation sur 500 installations réelles~:
$R^2 = 0{,}9$ sur la production mensuelle). Son idée centrale, que nous
reprenons, est de \textbf{séparer le problème en deux échelles}~:

\begin{enumerate}
  \item \textbf{La météo varie à l'échelle de quelques kilomètres.} Le
    rayonnement solaire incident (nuages, turbidité) est fourni par des données
    satellites sur une grille d'environ 4--5~km. Pour chaque cellule de cette
    grille, on pré-calcule une \emph{table}~: l'irradiance moyenne reçue par un
    plan incliné, pour toutes les combinaisons d'orientation et de pente de
    toit, pour chaque mois et chaque heure de la journée.
  \item \textbf{La géométrie varie à l'échelle de 50~cm.} La forme des toits
    (pente, orientation) et les ombres portées (relief, bâtiments voisins) sont
    extraites de modèles numériques d'élévation à 0,5~m issus du LiDAR. Pour
    chaque pixel de toiture, il suffit alors de \emph{lire} dans la table
    pré-calculée la valeur correspondant à sa pente et son orientation, et de
    supprimer le rayonnement direct aux heures où le soleil est caché par
    l'horizon local.
\end{enumerate}

Cette factorisation est ce qui rend la méthode scalable~: le calcul coûteux
(transposition du rayonnement, sections~\ref{sec:transposition}
et~\ref{sec:transpAgr}) est fait \textbf{une fois par cellule météo de
$0{,}05^\circ$} (environ 4~$\times$~5,5~km en France) au lieu d'une fois par
pixel de 0,5~m -- soit un facteur d'économie d'environ $10^8$.

\subsection{État d'avancement}

\begin{itemize}
  \item \textbf{Fait}~: détection des pixels de toiture exploitables (croisement
    LiDAR / BD TOPO), calcul de la pente et de l'orientation de chaque pixel,
    calcul des angles d'horizon, construction des tables d'irradiance par cellule
    météo (téléchargement PVGIS + transposition Perez), et la \textbf{phase
    tuile} (section~\ref{sec:tuile})~: croisement des deux branches, irradiation
    par pixel \textbf{avec ombrage par l'horizon}, agrégation par bâtiment,
    \textbf{estimation de la production PV} (modèle simple), et écriture du
    \textbf{GeoPackage de sortie}. La pipeline tourne de bout en bout pour une
    dalle et produit un fichier bâtiments exploitable.
  \item \textbf{À faire} (section~\ref{sec:afaire})~: étendre l'horizon aux dalles
    voisines (effet de bord), raffiner l'interpolation $(\alpha,\beta)$, le modèle
    PV complet, et le passage à l'échelle (boucle multi-dalles + fusion par
    région).
\end{itemize}

% ============================================================================
\section{Architecture générale}
% ============================================================================

\subsection{Les deux branches de la pipeline}

La pipeline travaille par \textbf{dalle} (tuile) de 1~km~$\times$~1~km, le
format de distribution des données LiDAR de l'IGN. Deux branches indépendantes
préparent les ingrédients, une phase finale les croise~:

\begin{center}
\begin{tabular}{ccc}
\textbf{Branche A « météo »} & & \textbf{Branche B « géométrie »}\\
(une fois par cellule de $0{,}05^\circ$) & & (une fois par dalle de 1~km)\\[2pt]
\fichier{irr\_dep.py} & & \fichier{main\_MNTMNS.py}\\
$\downarrow$ & & $\downarrow$\\
téléchargement PVGIS (19 ans horaires) & & emprise dalle + bâtiments BD TOPO\\
$\downarrow$ & & $\downarrow$\\
transposition Perez sur 192 plans & & masques toiture, pente, aspect\\
$\downarrow$ & & $\downarrow$\\
table \fichier{.npz} : $B, D, SAZ, SEL$ & & angles d'horizon (36 directions)\\[6pt]
\multicolumn{3}{c}{$\searrow$ \hspace{2cm} $\swarrow$}\\[2pt]
\multicolumn{3}{c}{\textbf{Phase tuile} (\fichier{irr\_calcul.py}, section~\ref{sec:tuile}) : irradiance par pixel + agrégation par bâtiment}\\
\multicolumn{3}{c}{$\downarrow$}\\
\multicolumn{3}{c}{\textbf{GeoPackage de sortie} : un bâtiment par ligne (irradiation, surfaces, orientation)}\\
\end{tabular}
\end{center}

\noindent L'angle d'horizon (branche B) est calculé puis \emph{utilisé} par la
phase tuile pour l'ombrage du rayonnement direct. La seule limite actuelle est
que l'horizon est calculé sur la dalle seule (sans ses voisines), donc
sous-estimé en bordure de dalle (section~\ref{sec:afaire}).

Les deux branches ne communiquent que par \textbf{des fichiers}~: des tables
\fichier{.npz} d'un côté, des GeoTIFF de l'autre. C'est volontaire~: chaque
étape peut être relancée, vérifiée ou parallélisée indépendamment, et un
plantage ne fait jamais perdre que l'étape en cours.

\subsection{Arborescence et exécution}

\begin{lstlisting}[language={}]
Evaluation_PV_toitures_France/
|-- main_MNTMNS.py            # orchestre UNE dalle (geometrie + irradiance), pour tests
|-- src/
|   |-- geometrie/            # branche geometrie (par dalle)
|   |   |-- ExtractBDtopo.py      # emprise de dalle + chargement BD TOPO
|   |   |-- clip_MNTMNS_pixel.py  # masques toiture + pente + aspect
|   |   `-- horizon.py            # angles d'horizon par raycast
|   `-- irradiance/
|       |-- irr_fct.py        # moteur lookup table (transposition + lecture)
|       |-- irr_dep.py        # campagne : construit la lookup table
|       |-- irr_calcul.py     # lookup table -> irradiance + production par batiment
|       `-- irr_tuile.py      # helper coords/centre d'une dalle (tests, hors pipeline finale)
|-- data/
|   |-- raw/                  # donnees brutes (MNS, MNT, BD TOPO .gpkg)
|   |-- tables/               # tables meteo .npz produites par irr_dep
|   `-- processed/            # sorties par dalle (masques, horizons, .gpkg)
`-- environment.yml           # environnement conda "stage-lidar"
\end{lstlisting}

Tout se lance \textbf{depuis la racine du projet}, dans l'environnement conda
\fichier{stage-lidar}~:
\begin{lstlisting}[language=bash]
python main_MNTMNS.py                 # une dalle complete (geometrie + irradiance)
python src/irradiance/irr_dep.py      # branche A (meteo, par bbox)
\end{lstlisting}

Les chemins relatifs (\fichier{data/...}) sont résolus par rapport au répertoire
courant, d'où l'importance de lancer depuis la racine.

% ============================================================================
\section{Les données d'entrée}
% ============================================================================

\subsection{LiDAR HD IGN : MNS et MNT}
\label{sec:lidar}

Le programme \textbf{LiDAR HD} de l'IGN couvre la France d'un nuage de points
laser aéroporté dense ($\geq 10$ points/m$^2$). L'IGN en dérive deux
\emph{modèles numériques d'élévation}, distribués en dalles GeoTIFF de
1~km~$\times$~1~km à 0,5~m de résolution (rasters de $2000 \times 2000$
pixels)~:

\begin{description}[style=nextline]
  \item[MNS (Modèle Numérique de Surface)] altitude du \emph{sursol}~: on y voit
    les toits, les arbres, tout ce qui dépasse. C'est lui qui porte la géométrie
    des toitures.
  \item[MNT (Modèle Numérique de Terrain)] altitude du \emph{sol nu}, comme si
    tous les objets étaient effacés.
\end{description}

La différence $\text{MNH} = \text{MNS} - \text{MNT}$ (modèle numérique de
\emph{hauteur}) donne la hauteur des objets au-dessus du sol~: $\approx 0$ sur
le sol, plusieurs mètres sur un toit ou un arbre. C'est le critère central de la
détection (section~\ref{sec:clip}).

\paragraph{Convention de nommage des dalles.} Un fichier s'appelle par exemple
\fichier{LHD\_FXX\_0495\_6611\_}\allowbreak\fichier{MNS\_O\_0M50\_}\allowbreak\fichier{LAMB93\_IGN69.tif}. Les deux nombres
\fichier{0495} et \fichier{6611} sont les coordonnées Lambert~93 \textbf{en km du
coin nord-ouest} de la dalle (vérifié sur les données~: pour la dalle
\fichier{0475\_6594}, l'emprise lue dans le GeoTIFF va de $y = 6\,593\,000$ à
$y = 6\,594\,000$~m). Attention, c'est une source d'erreur classique~: le nom
donne le coin \emph{haut}-gauche, pas le coin bas-gauche.

\paragraph{Valeur nodata.} Les pixels sans donnée valent $-9999$ dans les
fichiers IGN. Partout dans le code, ils sont convertis en \fonction{NaN}
(\emph{Not a Number}) dès la lecture, pour qu'aucun calcul ne les prenne pour
une altitude réelle de $-9999$~m.

\subsection{BD TOPO : les emprises de bâtiments}

La \textbf{BD TOPO} est la base vectorielle de référence de l'IGN. On utilise sa
couche \fichier{batiment}~: un polygone par bâtiment (l'emprise au sol, précision
métrique), avec des attributs. Elle est distribuée par département au format
\textbf{GeoPackage} (\fichier{.gpkg}), un conteneur SQLite standard pour données
géographiques. Les attributs conservés par la pipeline~:

\begin{center}
\begin{tabular}{ll}
\toprule
\fichier{cleabs} & identifiant unique IGN du bâtiment (clé de jointure finale)\\
\fichier{nature}, \fichier{usage\_1} & typologie (église, serre, industriel\ldots)\\
\fichier{hauteur}, \fichier{nombre\_d\_etages} & contrôles de cohérence possibles\\
\fichier{etat\_de\_l\_objet} & on ne garde que « En service »\\
\fichier{construction\_legere} & on exclut les structures légères (auvents\ldots)\\
\bottomrule
\end{tabular}
\end{center}

Le rôle de la BD TOPO est de dire \emph{où chercher des toits}~: sans elle, le
MNH seul ne distingue pas un toit d'un arbre.

\subsection{PVGIS et SARAH-3 : le rayonnement solaire}
\label{sec:pvgis}

\textbf{PVGIS} (Photovoltaic Geographical Information System) est un service web
gratuit du Centre commun de recherche de la Commission européenne (JRC). Son API
(version 5.3, \url{https://re.jrc.ec.europa.eu/api/v5_3/}) fournit des
séries temporelles \textbf{horaires} de rayonnement au sol, issues de la base
\textbf{SARAH-3}\footnote{SARAH-3 (CM SAF / EUMETSAT)~: rayonnement satellite à
$0{,}05^\circ$, pas 30~min, depuis 1983 -- voir Réf. PVGIS n'en utilise ici qu'un
sous-ensemble horaire 2005--2023.}~: des estimations dérivées des images du
satellite géostationnaire Météosat, sur une grille de
$0{,}05^\circ \times 0{,}05^\circ$ (environ $3{,}8 \times 5{,}5$~km à la latitude
de la France), disponibles de 2005 à 2023.

C'est l'équivalent (en plus récent) des données « CM SAF / Heliosat SARAH »
utilisées par Buffat. L'avantage du satellite sur les stations météo au sol~:
une couverture spatiale continue et homogène sur 19 ans, là où les stations
sont éparses et demanderaient une interpolation hasardeuse.

Pour chaque heure, PVGIS fournit (entre autres) les deux composantes du
rayonnement sur un plan horizontal -- direct et diffus -- définies en
section~\ref{sec:composantes}. La requête est faite avec une inclinaison nulle
précisément pour récupérer ces grandeurs horizontales de base.

\subsection{Systèmes de coordonnées}

Deux systèmes cohabitent, et la pipeline passe de l'un à l'autre explicitement~:

\begin{description}[style=nextline]
  \item[Lambert 93 (code EPSG 2154)] la projection plane officielle de la France
    métropolitaine. Coordonnées $(x, y)$ \textbf{en mètres}~; toutes les données
    IGN (LiDAR, BD TOPO) sont dans ce système. On peut y mesurer des distances
    et des surfaces directement.
  \item[WGS 84 (code EPSG 4326)] latitude/longitude \textbf{en degrés}, le
    système GPS. C'est ce qu'attend l'API PVGIS et ce qu'utilisent les calculs
    astronomiques de position du soleil.
\end{description}

La conversion est faite par la bibliothèque \fichier{pyproj}
(section~\ref{sec:irrtuile}). Détail important~: l'option
\fonction{always\_xy=True} impose l'ordre (longitude, latitude) -- sans elle,
pyproj suit l'ordre officiel du système, qui est (latitude, longitude) pour le
WGS84, source classique d'inversions silencieuses.

% ============================================================================
\section{Notions de rayonnement solaire}
\label{sec:notions}
% ============================================================================

Cette section pose les concepts physiques utilisés par la branche météo. Tout le
vocabulaire défini ici est réutilisé dans la description des fonctions.

\subsection{Les composantes du rayonnement}
\label{sec:composantes}

L'\textbf{irradiance} est une puissance surfacique instantanée, en W/m$^2$.
L'\textbf{irradiation} est son intégrale dans le temps, en Wh/m$^2$ (ou
kWh/m$^2$). Au sol, le rayonnement reçu par une surface se décompose en trois
parts d'origines différentes~:

\begin{description}[style=nextline]
  \item[Direct (\emph{beam})] les rayons qui arrivent en ligne droite du disque
    solaire. C'est la part qui crée les ombres, et donc la seule affectée par
    les masques (relief, bâtiments voisins).
  \item[Diffus] la lumière du ciel~: rayonnement diffusé par les molécules, les
    aérosols et les nuages, qui arrive de toute la voûte céleste. Par temps
    couvert, c'est la totalité du rayonnement.
  \item[Réfléchi] la part renvoyée par le sol environnant vers une surface
    inclinée. Elle est proportionnelle à l'\textbf{albédo} $\rho$ du sol (part
    réfléchie, sans dimension)~; on prend $\rho = 0{,}2$, valeur standard pour
    un environnement urbain/végétal moyen (même valeur que Buffat).
\end{description}

Sur un \textbf{plan horizontal}, on note (conventions internationales,
reprises par pvlib)~:
\begin{itemize}
  \item $GHI$ (\emph{Global Horizontal Irradiance})~: le total horizontal~;
  \item $BHI$ (\emph{Beam Horizontal Irradiance})~: la part directe~;
  \item $DHI$ (\emph{Diffuse Horizontal Irradiance})~: la part diffuse~;
\end{itemize}
avec la relation de fermeture $GHI = BHI + DHI$ (le réfléchi est nul sur un plan
horizontal, qui ne « voit » pas le sol).

Enfin, le \textbf{DNI} (\emph{Direct Normal Irradiance}) est l'irradiance
directe mesurée sur un plan \emph{perpendiculaire aux rayons}. C'est la grandeur
géométriquement « pure » du faisceau~: connaître le DNI permet de projeter le
direct sur n'importe quel plan. Lien avec le direct horizontal~:
\begin{equation}
  DNI = \frac{BHI}{\cos \theta_z},
  \label{eq:dni}
\end{equation}
où $\theta_z$ est l'angle zénithal du soleil (section~\ref{sec:soleil}).
Cette division devient numériquement explosive quand le soleil est à
l'horizon ($\cos\theta_z \to 0$)~: c'est pour cela que le code passe par une
fonction pvlib munie de garde-fous plutôt que par une division naïve
(section~\ref{sec:transpAgr}).

\subsection{La position du soleil}
\label{sec:soleil}

La position du soleil dans le ciel, vue d'un point donné à un instant donné, est
décrite par deux angles~:
\begin{itemize}
  \item l'\textbf{azimut} $\gamma_s$~: la direction horizontale, en convention
    boussole ($0^\circ$ = nord, $90^\circ$ = est, $180^\circ$ = sud,
    $270^\circ$ = ouest)~;
  \item l'\textbf{élévation} $\theta_h$ (hauteur au-dessus de l'horizon) ou,
    de façon équivalente, l'\textbf{angle zénithal} $\theta_z = 90^\circ - \theta_h$
    (écart à la verticale).
\end{itemize}

On utilise l'élévation \emph{apparente}~: l'atmosphère réfracte la lumière et
relève légèrement le soleil visible par rapport au soleil géométrique (environ
$0{,}5^\circ$ à l'horizon) -- c'est la bonne grandeur à comparer à un horizon
physique.

Le calcul est fait par l'algorithme \textbf{SPA} (\emph{Solar Position
Algorithm}, NREL), précis à $\pm 0{,}0003^\circ$, via
\fonction{pvlib.solarposition.get\_solarposition}. C'est exactement l'algorithme
employé par Buffat. Toutes les heures de la pipeline sont en \textbf{UTC}
(jamais en heure légale française)~: les séries PVGIS sont en UTC, et garder un
fuseau unique évite les pièges de l'heure d'été.

\subsection{La transposition : du plan horizontal au plan du toit}
\label{sec:transposition}

Les données météo donnent le rayonnement sur un plan \emph{horizontal}~; un toit
est un plan \emph{incliné}, caractérisé par sa pente $\beta$ (0$^\circ$ = plat,
90$^\circ$ = vertical) et son orientation (azimut) $\alpha$ (convention
boussole~: $\alpha = 180^\circ$ pour un toit plein sud). Passer de l'un à
l'autre s'appelle la \textbf{transposition}. Chaque composante se transpose
différemment (équations 1 à 4 de Buffat)~:

\paragraph{Le direct : pure géométrie.} Il suffit de projeter le faisceau sur le
plan~:
\begin{equation}
  G_{B,\alpha,\beta} \;=\; DNI \cdot \max(0, \cos\theta_i)
  \;=\; \frac{BHI}{\cos\theta_z}\,\cos\theta_i ,
  \label{eq:beam}
\end{equation}
où $\theta_i$ est l'\textbf{angle d'incidence}, l'angle entre les rayons et la
normale au plan~:
\begin{equation}
  \cos\theta_i = \cos\theta_z\cos\beta
              + \sin\theta_z\sin\beta\cos(\gamma_s - \alpha).
\end{equation}
Le $\max(0,\cdot)$ traduit le fait qu'un soleil derrière le plan ne l'éclaire pas.

\paragraph{Le diffus : le modèle de Perez.} Si le ciel était une voûte de
luminance uniforme (modèle « isotrope »), un plan de pente $\beta$ en verrait la
fraction $(1+\cos\beta)/2$. En réalité le ciel est plus brillant autour du
soleil (composante \emph{circumsolaire}) et près de l'horizon (\emph{bande
d'horizon}). Le modèle empirique de \textbf{Perez et al. (1990)} corrige ces
deux effets~:
\begin{equation}
  G_{D,\alpha,\beta} = DHI\left[(1-F_1)\,\frac{1+\cos\beta}{2}
      \;+\; F_1\,\frac{a}{b} \;+\; F_2\sin\beta\right],
  \label{eq:perez}
\end{equation}
où $F_1$ et $F_2$ sont des coefficients empiriques (poids du circumsolaire et de
la bande d'horizon) qui dépendent de l'état du ciel -- décrit par un indice de
clarté et une luminosité calculés à partir de $DHI$, $DNI$, de la masse d'air et
du rayonnement extraterrestre -- et $a/b$ est un terme géométrique de projection
du circumsolaire ($a = \max(0,\cos\theta_i)$, $b = \max(\cos 85^\circ,
\cos\theta_z)$). C'est le modèle retenu par Buffat (et par la communauté PV en
général) car il est nettement plus juste que l'isotrope, notamment pour des
plans verticaux ou orientés est/ouest. Il est implémenté dans pvlib~; le code ne
fait que l'appeler avec les bons ingrédients.

\paragraph{Le réfléchi : modèle isotrope au sol.} Le plan incliné « voit » le
sol avec le facteur de vue complémentaire $(1-\cos\beta)/2$, et le sol renvoie
$\rho \cdot GHI$~:
\begin{equation}
  G_{R,\alpha,\beta} = \rho \; GHI \; \frac{1-\cos\beta}{2}.
  \label{eq:reflechi}
\end{equation}

\paragraph{Total dans le plan.} La somme des trois (équation 4 de Buffat)~:
\begin{equation}
  G_{\alpha,\beta} = G_{B,\alpha,\beta} + G_{D,\alpha,\beta} + G_{R,\alpha,\beta}.
\end{equation}
Dans le code, le diffus ciel (Perez) et le réfléchi sont regroupés dans un même
tableau $D$ (« tout ce qui n'est pas du direct »), car c'est la part qui
\emph{subsiste} quand le soleil est masqué~; le direct $B$ est gardé séparé car
c'est lui qu'on annulera à l'ombre.

\subsection{L'ombrage par l'horizon}
\label{sec:ombrage}

Depuis un pixel de toit, dans une direction donnée, l'\textbf{angle d'horizon}
$\theta_{hor}(\varphi)$ est l'élévation angulaire du plus haut obstacle visible
(faîtage voisin, immeuble, colline). La règle d'ombrage de Buffat (équations 5
et 6) est binaire et ne s'applique qu'au direct~:
\begin{equation}
  \text{si } \theta_h \;\leq\; \theta_{hor}(\gamma_s)
  \quad\Longrightarrow\quad G_{B} = 0
  \quad (\text{le diffus et le réfléchi sont conservés}).
\end{equation}
Autrement dit~: à chaque pas de temps, on regarde l'azimut du soleil, on lit
l'angle d'horizon dans cette direction, et si le soleil est plus bas que
l'obstacle, le pixel est à l'ombre -- il ne reçoit que du diffus. C'est une
approximation (le diffus aussi est partiellement masqué par les obstacles), mais
c'est celle de l'étude de référence et elle capte l'essentiel du phénomène, le
direct représentant la majorité de l'énergie annuelle en France.

\subsection{De l'irradiance à l'irradiation : le profil (mois $\times$ heure)}
\label{sec:profil}

Plutôt que de garder 19 ans de séries horaires partout, la pipeline résume le
climat par un \textbf{profil moyen (mois, heure)}~: un tableau $12 \times 24$
dont la case $(m, h)$ contient la moyenne, sur les 19 années, de l'irradiance
aux heures $h$ UTC du mois $m$ (chaque case moyenne donc environ
$19 \times 30 \approx 570$ valeurs horaires). C'est l'équivalent du « jour type
mensuel » de Buffat (qui travaille à la demi-heure~; PVGIS fournit de l'horaire).

L'énergie s'en déduit exactement comme dans l'équation 7 de Buffat~: chaque case
vaut pour 1~heure du jour type, donc
\begin{equation}
  I_m \;=\; n_{\text{jours}}(m) \sum_{h=0}^{23} \bar G(m,h) \times 1\,\text{h}
  \qquad\text{[Wh/m}^2\text{/mois]},
  \qquad
  I_{\text{an}} = \frac{1}{1000}\sum_{m=1}^{12} I_m
  \quad\text{[kWh/m}^2\text{/an]}.
  \label{eq:energie}
\end{equation}

\paragraph{Pourquoi transposer \emph{avant} de moyenner ?} C'est un point
méthodologique important (et le cœur de \fonction{transpAgr}). La transposition
dépend de la position du soleil, qui change à chaque instant~; transposer une
moyenne reviendrait à éclairer le plan avec un « soleil moyen », ce qui fausse
les géométries (notamment matin/soir et hiver). On transpose donc \textbf{chaque
heure des 19 ans} avec sa vraie position de soleil, et on ne moyenne
qu'ensuite. C'est exactement l'ordre des opérations de Buffat
(« \emph{for each of the available half-hour records [\ldots] the irradiance
values on inclined surfaces are computed} », puis agrégation par mois).

% ============================================================================
\section{Branche B -- la géométrie : détection des toitures}
% ============================================================================

\subsection{\fichier{main\_MNTMNS.py} : l'orchestrateur}

Ce script enchaîne toute la chaîne pour \textbf{une dalle}, configurée en tête de
fichier (noms des fichiers MNS/MNT, chemin du GeoPackage BD TOPO, dossier de
sortie). Son déroulé~:

\begin{enumerate}
  \item vérification de l'existence des fichiers d'entrée (sortie immédiate avec
    message clair sinon)~;
  \item \fonction{tileBounds(MNS\_NAME)} $\rightarrow$ emprise de la dalle en
    Lambert 93~;
  \item \fonction{loadBuild(GPKG\_PATH, bounds)} $\rightarrow$ bâtiments BD TOPO
    de la dalle, filtrés~;
  \item \fonction{clip\_MNSMNT\_pixels(...)} $\rightarrow$ masques de pixels de
    toiture (inclinés / plats), pente et aspect par pixel~;
  \item \fonction{compHZ(...)} $\rightarrow$ angles d'horizon des pixels de
    toiture (36 directions de $10^\circ$), retournés en mémoire comme tableau
    $(N, 36)$ \emph{et} écrits en 36 rasters pour QGIS~;
  \item \fonction{irradianceTuile(...)} (phase tuile, section~\ref{sec:tuile})
    $\rightarrow$ irradiation par pixel \emph{avec ombrage}, agrégation par
    bâtiment, et écriture du GeoPackage de sortie
    \fichier{<MNS>\_irradiance.gpkg}.
\end{enumerate}

Chaque étape est chronométrée et affiche des compteurs -- utile pour estimer le
coût au moment de passer à des milliers de dalles. Le script ne contient aucune
logique métier~: tout est dans les modules de \fichier{src/}, ce qui permet de
réutiliser les fonctions ailleurs (notebooks, tests, futur traitement par lots).

\subsection{\fichier{src/geometrie/ExtractBDtopo.py}}

\subsubsection{\fonction{tileBounds(mns\_name)} : du nom de fichier à l'emprise}

\textbf{Rôle.} Déduire l'emprise géographique $(x_{\min}, y_{\min}, x_{\max},
y_{\max})$ d'une dalle, en mètres Lambert 93, à partir de son seul nom de
fichier -- sans ouvrir le raster.

\textbf{Fonctionnement.} Une expression régulière extrait les deux groupes de 4
chiffres encadrés de tirets bas~:

\begin{lstlisting}
match = re.search(r'_(\d{4})_(\d{4})_', mns_name)
x_km = int(match.group(1));  y_km = int(match.group(2))
x_min = x_km * 1000;  x_max = x_min + 1000
y_max = y_km * 1000;  y_min = y_max - 1000     # nom IGN = coin NORD-ouest
\end{lstlisting}

L'expression \fonction{r'\_(\textbackslash d\{4\})\_(\textbackslash d\{4\})\_'}
se lit~: « un tiret bas, puis exactement 4 chiffres (capturés), un tiret bas, 4
chiffres (capturés), un tiret bas ». Si le nom ne correspond pas, la fonction
lève une erreur explicite plutôt que de continuer avec des coordonnées fausses.
Le point délicat est la dernière ligne~: le nom encodant le coin
\emph{nord}-ouest (section~\ref{sec:lidar}), $y$ nommé est le \emph{haut} de la
dalle~; on retranche 1~km pour obtenir le bas.

\subsubsection{\fonction{loadBuild(gpkg\_path, tile\_bounds)} : charger et filtrer les bâtiments}

\textbf{Rôle.} Renvoyer un \fonction{GeoDataFrame} (le « tableau avec une
colonne géométrie » de la bibliothèque geopandas) contenant les bâtiments
exploitables de la dalle.

\textbf{Fonctionnement.}
\begin{lstlisting}
gdf = gpd.read_file(gpkg_path, layer='batiment', bbox=tile_bounds)
gdf = gdf[(gdf['etat_de_l_objet'] == 'En service') &
          (gdf['construction_legere'] == False)].copy().reset_index(drop=True)
gdf = gdf[['cleabs', 'nature', 'usage_1', 'hauteur',
           'nombre_d_etages', 'geometry']].copy()
\end{lstlisting}

Trois idées~:
\begin{itemize}
  \item \textbf{Lecture spatiale partielle.} L'argument \fonction{bbox} fait
    utiliser l'index spatial du GeoPackage~: seuls les bâtiments intersectant la
    dalle sont lus. Un département entier contient des centaines de milliers de
    bâtiments~; en lire $\sim$1000 prend une fraction de seconde au lieu de
    plusieurs dizaines.
  \item \textbf{Filtrage attributaire.} On écarte les bâtiments hors service
    (ruines, en projet) et les constructions légères, dont les « toits » ne
    portent pas de PV. Le filtre est volontairement minimal à ce stade~; un
    filtrage par \fichier{nature}/\fichier{usage\_1} (serres, monuments
    historiques\ldots) est en commentaire dans le code, à activer plus tard.
  \item \textbf{\fonction{reset\_index(drop=True)}.} Après filtrage, l'index du
    GeoDataFrame est renuméroté $0, 1, 2, \ldots$~: c'est cet index (décalé de
    $+1$, voir plus bas) qui servira d'identifiant de bâtiment dans les rasters,
    et il doit être dense et prévisible.
\end{itemize}

\subsection{\fichier{src/geometrie/clip\_MNTMNS\_pixel.py} : \fonction{clip\_MNSMNT\_pixels}}
\label{sec:clip}

\textbf{Rôle.} C'est la fonction de \textbf{détection}~: elle décide, pixel par
pixel (0,5~m), ce qui est une toiture exploitable, et calcule pour chacun sa
pente et son orientation. Elle prend les chemins du MNS et du MNT, le
GeoDataFrame des bâtiments, et deux seuils~; elle renvoie deux masques entiers
(toits inclinés / toits plats), la pente, l'aspect, et les métadonnées raster
nécessaires pour écrire les sorties.

\subsubsection{Lecture des rasters et notion de transform affine}

\begin{lstlisting}
with rasterio.open(mns_path) as src:
    mns        = src.read(1).astype(np.float32)  # bande 1 -> tableau numpy (2000, 2000)
    transform  = src.transform                   # lien indices <-> coordonnees terrain
    meta       = src.meta.copy()                 # dimensions, CRS, nodata, ...
    nodata     = src.nodata
    resolution = src.res[0]                      # 0.5 m
\end{lstlisting}

Un GeoTIFF est une matrice de pixels accompagnée d'une \textbf{transformation
affine} qui convertit les indices (colonne, ligne) en coordonnées terrain
$(x, y)$~: $x = a\cdot\text{col} + c$, $y = e\cdot\text{lig} + f$ avec $e < 0$
-- la ligne 0 est en \emph{haut} (au nord) et les lignes croissent vers le
\emph{sud}. Cette orientation « lignes vers le sud » est invisible tant qu'on
fait des calculs locaux, mais devient cruciale pour l'aspect et pour l'horizon
(sens des déplacements dans la matrice). Après lecture, les pixels nodata sont
convertis en NaN~:

\begin{lstlisting}
if nodata is not None:
    mns[mns == nodata] = np.nan
    mnt[mnt == nodata] = np.nan
\end{lstlisting}

\subsubsection{Le MNH et la rasterisation des bâtiments}

\begin{lstlisting}
mnh = mns - mnt                                  # hauteur au-dessus du sol

geometries = gdf.geometry.buffer(buffer)         # buffer = 1.2 m
masque_bat = rasterize(
    zip(geometries, gdf.index + 1),              # index + 1 : 0 = fond
    out_shape=mns.shape, transform=transform, fill=0, dtype="int32")
\end{lstlisting}

\begin{itemize}
  \item \textbf{Le buffer de 1,2~m} dilate chaque polygone~: les emprises
    BD TOPO sont les murs au sol, or les toits débordent (avant-toits), et il
    existe un léger décalage possible entre la BD et le LiDAR. Sans buffer, on
    perdrait systématiquement les bords de toits.
  \item \textbf{\fonction{rasterize}} « brûle » chaque polygone dans une grille
    alignée sur le MNS~: chaque pixel dont le centre tombe dans le polygone
    $i$ reçoit la valeur $i+1$, les autres restent à \fonction{fill=0}. On
    obtient d'un coup \emph{où} sont les bâtiments \emph{et lequel}~: le raster
    sert à la fois de masque et de carte d'identifiants, ce qui permettra
    d'agréger les résultats par bâtiment sans aucune boucle sur les polygones.
  \item \textbf{Pourquoi $+1$~?} L'index du GeoDataFrame commence à 0, qui est
    aussi la valeur de fond~: sans le décalage, le premier bâtiment de chaque
    dalle serait invisible (confondu avec le fond). Pour retrouver le bâtiment
    depuis un pixel~: \fonction{gdf.loc[valeur - 1]}.
\end{itemize}

\subsubsection{Le pré-masque de validité}

\begin{lstlisting}
pts_ok = ((masque_bat > 0) & (mnh >= mnh_min) &
          np.isfinite(mns) & np.isfinite(mnt))
\end{lstlisting}

Un pixel candidat doit (i)~être dans une emprise (buffée) de bâtiment, (ii)~être
à au moins \fonction{mnh\_min} $= 1{,}5$~m au-dessus du sol -- ce qui élimine le
sol nu, les terrasses basses, voitures et murets inclus dans le buffer -- et
(iii)~avoir des altitudes valides dans les deux modèles. Le seuil de 1,5~m est
un compromis~: assez haut pour exclure le bruit au sol, assez bas pour garder
les toits de plain-pied.

\subsubsection{Pente et aspect par différences finies}

\begin{lstlisting}
dz_dy, dz_dx = np.gradient(mns, resolution)
pente  = np.degrees(np.arctan(np.sqrt(dz_dx**2 + dz_dy**2)))
aspect = np.degrees(np.arctan2(-dz_dx, dz_dy)) % 360
\end{lstlisting}

\fonction{np.gradient} calcule les dérivées partielles de l'altitude par
différences finies \emph{centrées}~: pour chaque pixel, la variation d'altitude
entre ses deux voisins, divisée par la distance ($2 \times 0{,}5$~m). Le premier
tableau renvoyé est la dérivée le long de l'axe 0 (les lignes), le second le
long de l'axe 1 (les colonnes).

\begin{itemize}
  \item \textbf{La pente} est l'angle dont le plan local s'écarte de
    l'horizontale~: la norme du gradient $\|\nabla z\|$ est la montée par mètre
    parcouru, et $\beta = \arctan\|\nabla z\|$.
  \item \textbf{L'aspect} est la direction \emph{vers laquelle le toit
    descend} (la direction que « regarde » le pan), en convention boussole.
    Avec l'axe des lignes orienté vers le sud (cf.~transform), la direction de
    descente est $(-\partial z/\partial x_{\text{est}},\,
    -\partial z/\partial y_{\text{nord}}) = (-\fonction{dz\_dx},\,
    +\fonction{dz\_dy})$, et l'angle boussole d'un vecteur $(v_{\text{est}},
    v_{\text{nord}})$ est $\operatorname{atan2}(v_{\text{est}},
    v_{\text{nord}})$ -- d'où la formule. Le \fonction{\% 360} ramène le
    résultat de $[-180, 180]$ vers $[0, 360[$. Vérifié sur cas synthétique~: un
    plan montant vers l'est donne bien un aspect de $270^\circ$ (il regarde
    l'ouest).
\end{itemize}

\subsubsection{L'érosion morphologique : nettoyer les bords}

\begin{lstlisting}
kernel = np.ones((3, 3), dtype=bool)
pts_ok_erode = binary_erosion(pts_ok, structure=kernel)
\end{lstlisting}

L'\textbf{érosion binaire} (morphologie mathématique, via scipy) ne garde un
pixel \fonction{True} que si \emph{tout} son voisinage $3\times3$ est
\fonction{True}~: elle retire une couronne d'un pixel au bord du masque.
Pourquoi~? Parce que le gradient centré d'un pixel de bord de toit utilise un
voisin \emph{hors} du toit (la façade, le sol)~: sa pente et son aspect sont
contaminés (pentes de $80^\circ$ en bord de pan, aspects aléatoires). Plutôt que
de filtrer ces artefacts par des seuils fragiles, on supprime la source~: tout
pixel dont un voisin est invalide.

\subsubsection{Les masques finaux}

\begin{lstlisting}
masque_incline = np.where(pts_ok_erode & (pente >= 10) & (pente <= 70) &
                          (aspect >= 90) & (aspect <= 270), masque_bat, 0)
masque_plat    = np.where(pts_ok_erode & (pente < 10),      masque_bat, 0)
\end{lstlisting}

Deux familles, traitées différemment ensuite~:
\begin{itemize}
  \item \textbf{toits inclinés} ($10$--$70^\circ$)~: les panneaux suivent le
    pan~; la borne haute écarte les façades résiduelles et les artefacts. Le
    filtre $90 \leq \text{aspect} \leq 270$ ne conserve que les pans orientés
    de l'est à l'ouest en passant par le sud -- choix méthodologique assumé
    (les pans nord produisent peu en France), à réévaluer si l'on veut un
    cadastre exhaustif (section~\ref{sec:limites})~;
  \item \textbf{toits plats} ($< 10^\circ$)~: les panneaux y seraient posés sur
    châssis avec une inclinaison choisie~; leur potentiel se calculera avec un
    $(\alpha, \beta)$ conventionnel plutôt qu'avec la géométrie du pixel.
\end{itemize}
Les masques portent l'identifiant du bâtiment (\fonction{masque\_bat}, donc
index~$+1$) là où la condition est vraie, 0 ailleurs~: ce sont à la fois des
masques booléens et des cartes de correspondance pixel $\rightarrow$ bâtiment.

\subsubsection{Les exports de contrôle}

Si \fonction{debug\_dir} est fourni, la fonction écrit cinq GeoTIFF (MNH masqué,
pente filtrée, aspect, deux masques) compressés en LZW, à ouvrir dans QGIS pour
contrôle visuel. Deux détails de robustesse~: le dossier est créé si besoin, et
la conversion NaN $\rightarrow$ nodata n'est appliquée qu'aux rasters flottants
(\fonction{np.isnan} déclencherait une erreur sur les masques entiers, qui ne
peuvent pas contenir de NaN).

\subsection{\fichier{src/geometrie/horizon.py} : \fonction{compHZ}}
\label{sec:horizon}

\textbf{Rôle.} Pour chaque pixel de toiture et chaque direction azimutale (36
directions par pas de $10^\circ$), calculer l'angle d'horizon
$\theta_{hor}(\varphi)$ défini en section~\ref{sec:ombrage}~: l'élévation du
plus haut obstacle visible dans cette direction, à partir du MNS.
\textbf{Sortie principale}~: un tableau $(N, 36)$ en mémoire (un angle d'horizon
en degrés par pixel de toiture et par direction), \emph{retourné} pour être
utilisé directement par la phase tuile (section~\ref{sec:tuile}). En effet de
bord, la fonction écrit aussi un GeoTIFF par direction
(\fichier{horizon\_000.tif} \ldots \fichier{horizon\_350.tif}, $-9999$ hors
toiture, LZW) pour visualisation dans QGIS. L'ordre des lignes du tableau suit
\fonction{np.where(masque\_toiture)}~: c'est ce qui garantit que la ligne $k$ du
tableau horizon correspond au même pixel que la ligne $k$ côté phase tuile.

\subsubsection{Le principe : un lancer de rayons discret}

Pour un pixel d'altitude $z_0$ et une direction $\varphi$, on avance pas à pas
le long du rayon ($k = 1, 2, \ldots, k_{\max}$ pixels)~; à chaque pas, l'obstacle
potentiel situé à la distance $k\,\delta$ ($\delta = 0{,}5$~m) et à l'altitude
$z_k$ est vu sous l'angle
\begin{equation}
  \theta_k = \operatorname{atan2}\!\big(z_k - z_0,\; k\,\delta\big),
\end{equation}
et l'angle d'horizon est le maximum rencontré~:
$\theta_{hor}(\varphi) = \max_k \theta_k$. Avec
\fonction{max\_distance\_m} $= 100$~m (réglage actuel dans \fichier{main}),
$k_{\max} = 200$ pas. L'angle est initialisé à $-90^\circ$~; un environnement
entièrement plus bas que le pixel donne un horizon négatif, ce qui est correct
(un faîtage voit « sous » l'horizontale).

\subsubsection{La vectorisation : la clé des performances}

L'algorithme naïf ferait trois boucles imbriquées (pixels $\times$ directions
$\times$ pas)~: des milliards d'itérations Python, rédhibitoire. L'astuce est de
ne boucler en Python que sur les directions (36) et les pas (200), et de traiter
\textbf{tous les pixels de toiture d'un coup} à chaque pas, par indexation
numpy~:

\begin{lstlisting}
x, y = np.where(masque_toiture)        # lignes/colonnes des N pixels de toit
z0 = mns[x, y]                          # leurs N altitudes

for phi_deg in range(0, 360, step):
    dx = -np.cos(phi_rad)               # nord = lignes qui diminuent
    dy =  np.sin(phi_rad)               # est  = colonnes qui augmentent
    theta = np.full(len(x), -np.pi/2, dtype=np.float32)
    for k in range(1, max_dist_px + 1):
        x_k = np.clip(np.round(x + k*dx).astype(int), 0, H-1)
        y_k = np.clip(np.round(y + k*dy).astype(int), 0, W-1)
        angle = np.arctan2(mns[x_k, y_k] - z0, k * res)
        theta = np.fmax(theta, angle)   # max courant ; fmax ignore les NaN
\end{lstlisting}

Points à comprendre~:
\begin{itemize}
  \item \textbf{Conventions d'axes.} Dans la matrice, \fonction{x} désigne la
    \emph{ligne} et \fonction{y} la \emph{colonne} (ordre numpy). La ligne 0
    étant au nord, aller vers le nord c'est décrémenter la ligne, d'où
    $dx = -\cos\varphi$~; aller vers l'est c'est incrémenter la colonne, d'où
    $dy = +\sin\varphi$. Vérifié par test synthétique~: un mur de 10~m placé
    15~m au nord d'un pixel donne $\theta_{hor}(0^\circ) =
    \arctan(10/15) = 33{,}69^\circ$, exactement la valeur attendue.
  \item \textbf{Échantillonnage au plus proche voisin.} \fonction{np.round}
    arrondit la position réelle du rayon au pixel le plus proche -- une forme
    discrète de tracé de ligne, suffisante à 0,5~m de résolution.
  \item \textbf{\fonction{np.clip} aux bords.} Quand le rayon sort de la dalle,
    on ré-échantillonne le pixel de bord~: l'angle correspondant décroît avec la
    distance et n'affecte plus le maximum. Conséquence assumée~: près des bords
    de dalle, l'horizon est sous-estimé (les obstacles de la dalle voisine sont
    invisibles)~; la correction (charger un voisinage) est listée dans le
    travail restant.
  \item \textbf{\fonction{np.fmax} et non \fonction{np.maximum}.} Les deux
    calculent le maximum élément par élément, mais \fonction{np.maximum}
    \emph{propage} les NaN~: un seul pixel nodata sur le rayon aurait
    définitivement contaminé l'angle. \fonction{np.fmax} ignore les NaN, ce qui
    revient à traiter les trous de données comme transparents.
  \item \textbf{Coût.} $36 \times 200$ opérations vectorisées sur $N$ pixels de
    toiture~: quelques secondes par dalle pour $N \sim 10^5$. La distance de
    recherche est le paramètre de coût principal ($100$~m actuellement~;
    Buffat utilise 500~m pour l'horizon proche, complété d'un horizon lointain
    -- voir section~\ref{sec:afaire}).
\end{itemize}

Pourquoi \textbf{36 directions}~? L'horizon varie lentement avec l'azimut~; au
moment du masquage, on lira la direction la plus proche de l'azimut solaire
(erreur maximale $5^\circ$, du même ordre que le pas de $5^\circ$ utilisé par
Buffat). Et pourquoi \textbf{un raster par direction}~? Pour rester en 2D simple
(facile à ouvrir dans QGIS), au prix de 36 fichiers -- compressés LZW, ils sont
essentiellement vides hors toiture et pèsent peu.

% ============================================================================
\section{Branche A -- la météo : tables d'irradiance par cellule}
% ============================================================================

\subsection{\fichier{src/irradiance/irr\_fct.py} : le module central}

Ce module contient les constantes et toutes les fonctions réutilisables de la
branche météo. Il est conçu pour être importable depuis n'importe où
(\fonction{from src.irradiance.irr\_fct import *} depuis la racine).

\subsubsection{Les constantes}

\begin{lstlisting}
ALBEDO  = 0.2                     # albedo du sol (Buffat 2018)
ALPHAS  = np.arange(0, 360, 15)   # 24 orientations (0=Nord, 90=Est, ...)
BETAS   = np.arange(0, 71, 10)    # 8 pentes, 0 a 70 degres
N_JOURS = np.array([31, 28, 31, 30, 31, 30, 31, 31, 30, 31, 30, 31])
PAS     = 0.05                    # pas de la grille meteo PVGIS (degres)
DOSSIER = "data/tables"           # un .npz par cellule meteo
URL     = "https://re.jrc.ec.europa.eu/api/v5_3/"   # PVGIS 5.3 (SARAH-3)
\end{lstlisting}

La grille $(\alpha, \beta)$ est volontairement \textbf{grossière}
($15^\circ \times 10^\circ$, 192 combinaisons)~: l'irradiance annuelle varie
lentement et régulièrement avec l'orientation et la pente (c'est une surface
lisse, cf.~figure~2 de Buffat), donc une interpolation depuis cette grille
suffira au niveau du pixel. Buffat utilise exactement ces pas ($15^\circ$ en
orientation, $10^\circ$ en pente). La borne de $70^\circ$ couvre la gamme
acceptée par le masque incliné ($10$--$70^\circ$)~: aucune extrapolation ne sera
nécessaire.

\subsubsection{\fonction{telecharger(lat, lon)} : la requête PVGIS}

\textbf{Rôle.} Récupérer la série horaire 2005--2023 des composantes du
rayonnement horizontal pour la cellule SARAH-3 contenant le point demandé.

\begin{lstlisting}
out = pvlib.iotools.get_pvgis_hourly(
    lat, lon, start=2005, end=2023,
    raddatabase="PVGIS-SARAH3",
    components=True, surface_tilt=0, surface_azimuth=0,
    usehorizon=False,
    url=URL, map_variables=True)
df = out[0]
\end{lstlisting}

Décodage des arguments~:
\begin{itemize}
  \item \fonction{components=True} demande les composantes séparées plutôt que
    le seul total~; \fonction{surface\_tilt=0} demande le plan
    \emph{horizontal}, si bien que la composante « directe dans le plan »
    renvoyée (\fonction{poa\_direct}) \textbf{est} le $BHI$, et la « diffuse
    ciel » (\fonction{poa\_sky\_diffuse}) \textbf{est} le $DHI$
    (section~\ref{sec:composantes})~; la part réfléchie est nulle à plat.
  \item \fonction{usehorizon=False}~: PVGIS sait appliquer l'ombrage du relief
    lointain à ses séries, mais nous gérons l'horizon nous-mêmes (et à bien plus
    fine échelle)~; on demande donc le rayonnement \emph{non masqué} pour ne pas
    compter l'ombrage deux fois. Garder l'horizon lointain PVGIS de côté reste
    une option notée dans le code.
  \item \fonction{map\_variables=True} renomme les colonnes PVGIS
    (\fichier{Gb(i)}, \fichier{Gd(i)}\ldots) en noms standard pvlib.
  \item Le résultat \fonction{out} est un tuple dont le premier élément est le
    DataFrame (indexé par des horodatages UTC)~: $19$ ans $\times$ $8760$~h
    $\approx 166\,500$ lignes.
\end{itemize}

\subsubsection{\fonction{transpAgr(bhi, dhi, lat, lon)} : le cœur du calcul}
\label{sec:transpAgr}

\textbf{Rôle.} Transformer 19 ans de séries horaires \emph{horizontales} en
\textbf{tables de profils moyens sur plans inclinés}~: pour chacune des
$24 \times 8$ combinaisons $(\alpha, \beta)$, le profil $12\times24$ (mois
$\times$ heure UTC) de l'irradiance directe $B$ et diffuse-plus-réfléchie $D$~;
plus les profils $12\times24$ de position moyenne du soleil ($SAZ$~: azimut,
$SEL$~: élévation apparente). C'est l'implémentation de la phase~1 de Buffat
(sections \ref{sec:transposition} et \ref{sec:profil}).

\textbf{Étape 1 -- préparation des ingrédients} (une seule fois, hors boucle)~:

\begin{lstlisting}
times = bhi.index
ghi = (bhi + dhi).clip(lower=0)                       # relation de fermeture

sp = pvlib.solarposition.get_solarposition(times, lat, lon)   # SPA, ~166 500 instants
dni = pvlib.irradiance.dni(ghi, dhi, sp["apparent_zenith"]).fillna(0)
dni_extra = pvlib.irradiance.get_extra_radiation(times)
airmass = pvlib.atmosphere.get_relative_airmass(sp["apparent_zenith"])
cles = [times.month, times.hour]
\end{lstlisting}

\begin{itemize}
  \item \fonction{sp} contient pour chaque heure le zénith apparent, l'azimut,
    l'élévation apparente du soleil (algorithme SPA,
    section~\ref{sec:soleil}), calculés au centre de la cellule météo -- la
    position du soleil varie de façon négligeable à l'échelle d'une cellule de
    5~km.
  \item \fonction{pvlib.irradiance.dni} reconstruit le $DNI$ par
    l'équation~\eqref{eq:dni}, \emph{avec garde-fous}~: les valeurs négatives
    (bruit de mesure à l'aube) et les valeurs au-delà de $88^\circ$ de zénith
    (division quasi nulle, donc instable) sont neutralisées (NaN puis 0 via
    \fonction{fillna})~; c'est la raison pour laquelle on n'écrit pas la
    division à la main.
  \item \fonction{dni\_extra} est le rayonnement hors atmosphère
    ($\approx 1361$~W/m$^2$, modulé de $\pm 3{,}3\,\%$ par l'excentricité de
    l'orbite terrestre) et \fonction{airmass} la masse d'air relative
    (l'épaisseur d'atmosphère traversée, 1 au zénith)~: deux ingrédients
    requis par le modèle de Perez, calculés une seule fois et passés
    explicitement pour éviter 192 recalculs.
  \item \fonction{cles} prépare les clés de groupement (mois, heure) -- voir
    \fonction{profMH}.
\end{itemize}

\textbf{Étape 2 -- transposition sur les 192 plans}~:

\begin{lstlisting}
B = np.zeros((len(ALPHAS), len(BETAS), 12, 24), np.float32)
D = np.zeros_like(B)
for i, a in enumerate(ALPHAS):
    for j, b in enumerate(BETAS):
        poa = pvlib.irradiance.get_total_irradiance(
            surface_tilt=b, surface_azimuth=a,
            solar_zenith=sp["apparent_zenith"], solar_azimuth=sp["azimuth"],
            dni=dni, ghi=ghi, dhi=dhi,
            dni_extra=dni_extra, airmass=airmass, albedo=ALBEDO, model="perez")
        direct = poa["poa_direct"].fillna(0)
        diffus = (poa["poa_sky_diffuse"] + poa["poa_ground_diffuse"]).fillna(0)
        B[i, j] = profMH(direct, cles)
        D[i, j] = profMH(diffus, cles)
\end{lstlisting}

\fonction{get\_total\_irradiance} est la fonction pvlib qui applique d'un coup
les équations \eqref{eq:beam}, \eqref{eq:perez} et \eqref{eq:reflechi} à toute
la série temporelle (calcul vectorisé~: chaque appel traite les
166\,500 heures d'un coup~; la boucle Python ne fait que 192 tours). Elle
renvoie les composantes séparées dans le plan~; on regroupe diffus ciel
$+$ réfléchi sol dans $D$ (la part insensible au masquage) et on garde le direct
seul dans $B$ (la part qui sera annulée à l'ombre,
section~\ref{sec:ombrage}). Les \fonction{fillna(0)} neutralisent les NaN que
Perez produit la nuit (masse d'air non définie sous l'horizon) -- de
l'irradiance nulle, donc.

Chaque profil est immédiatement réduit en moyenne (mois, heure) par
\fonction{profMH}~: on ne stocke jamais $192$ séries de 166\,500 points, mais
$192$ tableaux $12\times24$. La mémoire reste constante.

\textbf{Étape 3 -- profils de position solaire}~:

\begin{lstlisting}
SAZ = profMH(sp["azimuth"], cles)
SEL = profMH(sp["apparent_elevation"], cles)
return B, D, SAZ, SEL
\end{lstlisting}

Ces deux tableaux $12\times24$ donnent la position \emph{moyenne} du soleil pour
chaque case du jour type -- exactement ce qu'il faudra, en phase tuile, pour
comparer le soleil à l'horizon de chaque pixel (équation 6 de Buffat) \emph{sans
jamais relancer de calcul astronomique}. Moyenner un azimut est en général
risqué (discontinuité $360 \to 0$), mais sans danger ici~: en France
métropolitaine, le soleil ne passe jamais au nord pendant la journée~; la
discontinuité ne concerne que des heures nocturnes, jamais utilisées (le direct
y est nul).

\paragraph{Coût mesuré.} Environ une minute par cellule, téléchargement compris
(192 transpositions de 19 ans). C'est le prix payé \emph{une fois par cellule de
$0{,}05^\circ$} -- ensuite, tous les pixels de toutes les dalles de la cellule
se servent dans la table.

\subsubsection{\fonction{profMH(serie, cles)} : la réduction (mois, heure)}

\textbf{Rôle.} Transformer une série temporelle quelconque en son profil moyen
$12 \times 24$ (section~\ref{sec:profil}).

\begin{lstlisting}
g = serie.groupby(cles).mean()          # cles = [times.month, times.hour]
out = np.zeros((12, 24), np.float32)
for (m, h), v in g.items():
    out[int(m) - 1, int(h)] = v
return out
\end{lstlisting}

Le \fonction{groupby} de pandas range chaque valeur horaire dans sa case (mois,
heure) et calcule la moyenne de chaque case~: une opération standard et très
optimisée. La petite boucle finale ($\leq 288$ tours) ne fait que recopier le
résultat dans un tableau numpy dense, en décalant le mois (1--12) vers l'indice
(0--11). Les cases sans données resteraient à 0 -- cas qui ne se produit pas
avec des séries PVGIS complètes, mais la fonction est robuste à des trous.

\subsubsection{\fonction{totAb(B, D)} : le contrôle annuel}
\label{sec:totAb}

\textbf{Rôle.} Réduire des tables $(\alpha, \beta, 12, 24)$ en irradiation
annuelle par $(\alpha, \beta)$, en kWh/m$^2$/an -- l'équation~\eqref{eq:energie}
sommée sur l'année~:

\begin{lstlisting}
E = (B + D).astype(np.float64)                        # W/m2 moyens par case
return (E * N_JOURS[None, None, :, None]).sum(axis=(2, 3)) / 1000.0
\end{lstlisting}

Chaque case (mois, heure) est une puissance moyenne qui s'applique 1~h par jour
type~: multipliée par le nombre de jours du mois (diffusion numpy sur l'axe
mois~: \fonction{N\_JOURS[None,None,:,None]} aligne le vecteur de 12 jours sur
le bon axe du tableau 4D), sommée sur mois et heures, divisée par 1000
(Wh $\to$ kWh). Le passage en \fonction{float64} évite les pertes de précision
dans la somme de $288$ termes stockés en \fonction{float32}.

C'est l'outil de \textbf{validation}~: sur la cellule (46{,}30~N, 0{,}05~E),
il donne 1346~kWh/m$^2$/an à plat, 1582 plein sud à $30^\circ$, 915 plein nord à
$30^\circ$ -- valeurs conformes à ce que PVGIS affiche pour la région
(Poitou). Servira aussi en phase tuile pour produire la carte annuelle.

\subsubsection{\fonction{cheminTable} et \fonction{chargerTable} : le stockage}

\textbf{Rôle.} Définir où vivent les tables et comment les relire.

\begin{lstlisting}
def cheminTable(lat, lon):
    return os.path.join(DOSSIER, f"table_{lat + 0.0:.2f}_{lon + 0.0:.2f}.npz")

_cache = {}
def chargerTable(lat, lon):
    la = round(round(lat / PAS) * PAS, 2)     # centre de cellule le plus proche
    lo = round(round(lon / PAS) * PAS, 2)
    if (la, lo) not in _cache:
        d = np.load(cheminTable(la, lo))
        _cache[(la, lo)] = (d["B"], d["D"], d["SAZ"], d["SEL"])
    return _cache[(la, lo)]
\end{lstlisting}

\begin{itemize}
  \item Le nom encode le centre de cellule arrondi à deux décimales
    (\fichier{table\_46.30\_0.05.npz}). Le \fonction{+ 0.0} est une subtilité
    flottante~: près du méridien de Greenwich (qui traverse la France), un
    arrondi peut produire le « zéro négatif » \fonction{-0.0}, qui s'afficherait
    \fichier{-0.00} et créerait deux noms pour la même cellule~; lui ajouter
    $0{,}0$ force le zéro positif.
  \item \fonction{chargerTable} accepte \emph{n'importe quel} point (typiquement
    le centre d'une dalle LiDAR) et retrouve la cellule météo qui le contient
    par arrondi au multiple de $0{,}05^\circ$. C'est \textbf{le point d'entrée}
    de la future phase tuile.
  \item Le cache (un dictionnaire de module) garantit une seule lecture disque
    par cellule et par exécution -- des dalles voisines partagent la même
    cellule météo.
\end{itemize}

\subsection{\fichier{src/irradiance/irr\_dep.py} : le script de campagne}

Ce script construit les tables pour une zone d'étude. La zone est une bbox
WGS84 configurée en tête de fichier (actuellement un rectangle sur la Vienne,
$7 \times 8 = 56$ cellules).

\subsubsection{\fonction{grilleCellules()}}

Génère la liste des centres de cellules couvrant la bbox~:
\fonction{np.arange} de \fonction{LAT\_MIN} à \fonction{LAT\_MAX} par pas de
$0{,}05$, idem en longitude, produit cartésien. Le \fonction{np.round(..., 2)}
corrige les erreurs d'accumulation flottante d'\fonction{arange} (qui peut
produire $46{,}300000000004$), pour que les noms de fichiers soient exacts.

\subsubsection{\fonction{construireCellule(lat, lon)}}

L'unité de travail~: une cellule de bout en bout.

\begin{lstlisting}
df = telecharger(lat, lon)
B, D, SAZ, SEL = transpAgr(df["poa_direct"], df["poa_sky_diffuse"], lat, lon)
np.savez_compressed(cheminTable(lat, lon), B=B, D=D, SAZ=SAZ, SEL=SEL,
                    alphas=ALPHAS, betas=BETAS, lat=lat, lon=lon,
                    source="PVGIS-SARAH3 2005-2023, Perez")
\end{lstlisting}

Le fichier \fichier{.npz} (archive numpy compressée) contient les quatre
tableaux \emph{et} leurs axes (\fonction{alphas}, \fonction{betas}) et
métadonnées (position, source)~: chaque table est auto-décrite, on peut la
relire dans deux ans sans ce document. Taille~: $\approx 155$~ko par cellule
(les tableaux $B$ et $D$ font $24 \times 8 \times 12 \times 24$ valeurs
\fonction{float32}, soit 0,44~Mo bruts, bien compressés car pleins de zéros
nocturnes). La France entière ($\sim 43\,000$ cellules) tiendrait dans
$\sim 7$~Go~: le pré-calcul national est réaliste.

\subsubsection{\fonction{main()} : la boucle robuste}

\begin{lstlisting}
for k, (lat, lon) in enumerate(pts, 1):
    if os.path.exists(cheminTable(lat, lon)):
        continue                          # deja faite : reprise gratuite
    try:
        construireCellule(lat, lon)
    except Exception as e:
        print(f"... ECHEC : {e}"); ratees.append((lat, lon))
    time.sleep(PAUSE)                     # politesse API (1 s)
\end{lstlisting}

Trois mécanismes pour une campagne de téléchargement fiable~:
\textbf{(i)} l'idempotence -- une cellule déjà construite est sautée, on peut
interrompre et relancer à volonté~; \textbf{(ii)} l'isolation des échecs -- une
erreur (réseau, API) est affichée et mémorisée mais n'arrête pas la campagne, et
une simple relance retentera les manquantes~; \textbf{(iii)} la pause de 1~s
entre requêtes, par courtoisie envers l'API publique (limite officielle~:
30~requêtes/s). Coût mesuré~: $\approx 1$~minute par cellule, dominé par le
calcul de transposition.

\subsection{\fichier{src/irradiance/irr\_tuile.py} : le pont dalle $\to$ cellule (helper)}
\label{sec:irrtuile}

Deux petites fonctions qui font le lien entre le monde IGN (dalles Lambert 93)
et le monde météo (cellules WGS84). C'est un \textbf{helper pour les tests} sur
une dalle isolée~; il ne fait pas partie de la pipeline finale (qui bouclera sur
les dalles avec sa propre logique de nommage)~:

\begin{lstlisting}
def nomCoord(MNS_name):              # "LHD_FXX_0495_6611_..." -> (495, 6611)
    p = MNS_name.replace(".tif", "").split("_")
    return int(p[2]), int(p[3])      # coin NORD-ouest, en km Lambert 93

def centreWGS84(x_km, y_km):
    xc, yc = x_km * 1000 + 500, y_km * 1000 - 500   # -500 : le nom = coin NORD-ouest
    tr = Transformer.from_crs(2154, 4326, always_xy=True)
    lon, lat = tr.transform(xc, yc)
    return lat, lon
\end{lstlisting}

\fonction{centreWGS84} place le centre de la dalle ($+500$~m vers l'est,
$-500$~m vers le sud depuis le coin nord-ouest) puis le convertit en
latitude/longitude avec pyproj (vérifié contre l'emprise réelle d'une dalle~:
écart $< 1$~m). C'est ainsi que la phase tuile retrouve la cellule météo de la
dalle~: \fonction{chargerTable(*centreWGS84(*nomCoord(MNS\_NAME)))}.
Une dalle de 1~km étant bien plus petite qu'une cellule météo de
$\sim 4 \times 5{,}5$~km, prendre la cellule du centre de dalle pour toute la
dalle est une approximation tout à fait acceptable.

\medskip\noindent
Le fichier \fichier{irr\_calcul.py} (section~\ref{sec:tuile}) contient, lui, les
fonctions de calcul réutilisées par la pipeline complète~: \fonction{irrPixels},
\fonction{masquerHorizon}, \fonction{agregerBatiment} et \fonction{irradianceTuile}.

% ============================================================================
\section{La phase tuile : du pixel au GeoPackage}
\label{sec:tuile}
% ============================================================================

C'est l'étape qui \textbf{croise les deux branches} et produit le fichier de
sortie. Elle vit dans \fichier{irr\_calcul.py} (le cœur de calcul réutilisable
par la pipeline complète) et est appelée en dernier par \fichier{main\_MNTMNS.py}.
Elle applique l'\textbf{ombrage par l'horizon}~; le seul raffinement reporté est
l'horizon des dalles voisines (section~\ref{sec:afaire}).

\subsection{L'idée : un profil par pixel, masqué, puis intégré}

L'ombrage agit \emph{heure par heure} (le soleil n'est caché qu'à certains
moments de la journée), donc un total annuel par orientation ne suffit pas~: il
faut le \textbf{profil $(12\times24)$} de chaque pixel. La phase tuile suit donc
exactement la phase~3 de Buffat~:
\begin{enumerate}
  \item pour chaque pixel, lire le profil direct $B$ et diffus $D$ de la case
    $(\alpha,\beta)$ la plus proche dans les tables météo~;
  \item \textbf{éteindre le direct} aux cases (mois, heure) où le soleil est sous
    l'horizon du pixel (\fonction{masquerHorizon})~;
  \item \textbf{intégrer} $B_{\text{masqué}} + D$ sur l'année (kWh/m$^2$/an), puis
    multiplier par la surface réelle.
\end{enumerate}
Un profil pèse $288$ valeurs~; pour les $N$ pixels d'une dalle on manipule donc
des tableaux $(N, 12, 24)$ -- quelques centaines de Mo, sans problème.

\subsection{\fonction{irrPixels} : profil, ombrage et énergie de chaque pixel}

\textbf{Rôle.} Construire un tableau avec \textbf{une ligne par pixel de
toiture}, portant son irradiation annuelle (ombrage compris), sa surface réelle,
et de quoi l'agréger ensuite.

\begin{lstlisting}
masque_bat = np.where(masque_incline > 0, masque_incline, masque_plat)
toit = masque_bat > 0
p = pente[toit];  a = aspect[toit]

i = np.round(a / pas_a).astype(int) % len(alphas)   # plus proche orientation
j = np.clip(np.round(p / pas_b).astype(int), 0, len(betas) - 1)  # plus proche pente
B_pix = B[i, j]                       # (N, 12, 24) direct
D_pix = D[i, j]                       # (N, 12, 24) diffus + reflechi

B_pix = masquerHorizon(B_pix, horizon, SAZ, SEL)    # ombrage du direct

e_m2 = ((B_pix + D_pix) * N_JOURS[None, :, None]).sum(axis=(1, 2)) / 1000.0
surf = 0.25 / np.cos(np.radians(p))                 # surface reelle du pixel (m2)

return pd.DataFrame({
    "id":      masque_bat[toit] - 1,                # index du batiment dans gdf
    "energie": e_m2 * surf,                         # kWh/an
    "surf":    surf,
    "pente":   p,
    "secteur": np.round(a / 45).astype(int) % 8,    # secteur d'orientation (0=N)
    "incline": masque_incline[toit] > 0,
})
\end{lstlisting}

Points à comprendre~:
\begin{itemize}
  \item \textbf{Fusion des deux masques.} \fonction{masque\_bat} réunit toits
    inclinés et plats~: chaque pixel de toit porte l'identifiant ($+1$) de son
    bâtiment, $0$ ailleurs. Les deux masques sont exclusifs (seuil de pente),
    donc aucun conflit. L'ordre des pixels de \fonction{toit} est identique à
    celui du tableau \fonction{horizon} (même masque, même parcours en lignes),
    ce qui aligne les deux -- une assertion le vérifie.
  \item \textbf{Lecture « plus proche voisin ».} Plutôt qu'une interpolation, on
    arrondit l'aspect au pas de la grille ($15^\circ$) et la pente au sien
    ($10^\circ$) pour tomber sur la case la plus proche des tables. Le
    \fonction{\% len(alphas)} gère naturellement le bouclage de l'azimut
    ($357^\circ \to$ case $0^\circ$). \fonction{B[i,j]} extrait alors le profil
    $(12,24)$ des $N$ pixels d'un coup. Une interpolation bilinéaire (plus douce)
    est un raffinement futur (section~\ref{sec:afaire}).
  \item \textbf{Intégration (= \fonction{totAb} par pixel).} Chaque case (mois,
    heure) est une puissance moyenne valant 1~h par jour type~: on multiplie par
    le nombre de jours du mois, on somme sur l'année, on divise par 1000
    (Wh~$\to$~kWh). C'est l'équation~\eqref{eq:energie} appliquée pixel par pixel,
    \emph{après} masquage du direct.
  \item \textbf{Surface réelle vs projetée.} Un pixel couvre $0{,}5\times0{,}5 =
    0{,}25$~m$^2$ \emph{au sol}~; la surface réelle du pan incliné est
    $0{,}25/\cos\beta$ (un toit à $45^\circ$ a $\sqrt2$ fois plus de surface que
    son ombre au sol). C'est elle qui porte les panneaux.
  \item \textbf{Le secteur d'orientation.} $\operatorname{round}(a/45)\bmod 8$
    range l'aspect dans l'un des 8 secteurs de $45^\circ$ ($0=$N, $2=$E, $4=$S,
    $6=$O)~: base de la « surface par orientation » (voir plus bas).
  \item Tout est \textbf{vectorisé}~: aucune boucle sur les pixels.
\end{itemize}

\subsection{\fonction{masquerHorizon} : l'ombrage du direct}

\textbf{Rôle.} Mettre à zéro le rayonnement direct des cases (mois, heure) où le
soleil est caché par l'horizon local du pixel -- la règle de Buffat
(équation~6, section~\ref{sec:ombrage}). Le diffus n'est jamais touché.

\begin{lstlisting}
n_dir = horizon.shape[1]
d = np.round(SAZ / (360 / n_dir)).astype(int) % n_dir   # direction ~ azimut du soleil
horizon_soleil = horizon[:, d]                           # (N, 12, 24)
B_pix[SEL[None, :, :] <= horizon_soleil] = 0.0           # soleil sous l'horizon -> ombre
\end{lstlisting}

\begin{itemize}
  \item \textbf{Quelle direction d'horizon ?} Pour chaque case (mois, heure), le
    soleil est à l'azimut $SAZ[m,h]$~; on prend la direction d'horizon la plus
    proche ($SAZ$ arrondi au pas de $10^\circ$, modulo 36). \fonction{horizon[:,
    d]} avec \fonction{d} de forme $(12,24)$ produit, par indexation avancée, le
    tableau $(N,12,24)$ donnant pour chaque pixel l'horizon dans la direction du
    soleil à chaque instant.
  \item \textbf{Le test d'ombre.} Si l'élévation du soleil $SEL[m,h]$ est
    inférieure ou égale à cet horizon, le pixel est à l'ombre à cet instant~: son
    direct passe à $0$ (le diffus et le réfléchi arrivent toujours).
  \item \textbf{La nuit est neutre.} Aux heures nocturnes le direct est déjà nul
    (rien à éteindre)~; le masquage n'a donc aucun effet parasite, même si
    l'azimut moyen de nuit (proche du nord) est mal défini.
  \item \textbf{Sur place.} On modifie \fonction{B\_pix} directement (c'est déjà
    une copie issue de \fonction{B[i,j]}), pour éviter d'allouer un second
    tableau $(N,12,24)$.
\end{itemize}

\subsection{\fonction{agregerBatiment} : sommer par bâtiment}

\textbf{Rôle.} Regrouper les pixels par bâtiment (« zonal statistics ») et
joindre le résultat aux polygones BD TOPO.

\begin{lstlisting}
inc = df[df.incline]                          # pixels inclines seulement
res = pd.DataFrame({
    "irr_an_kwh":       df.groupby("id").energie.sum(),
    "surf_inclinee_m2": inc.groupby("id").surf.sum(),
    "surf_plate_m2":    df[~df.incline].groupby("id").surf.sum(),
    "pente_moy":        inc.groupby("id").pente.mean(),
    "nb_pixels":        df.groupby("id").size(),
})
for s, nom in enumerate(SECTEURS):            # surface inclinee par orientation
    res[f"surf_{nom}_m2"] = inc[inc.secteur == s].groupby("id").surf.sum()

res = res.fillna(0.0)
res["surf_tot_m2"]   = res.surf_inclinee_m2 + res.surf_plate_m2
res["irr_an_kwh_m2"] = res.irr_an_kwh / res.surf_tot_m2
return gdf.join(res, how="inner")             # batiments avec >= 1 pixel de toit
\end{lstlisting}

\begin{itemize}
  \item \textbf{\fonction{groupby("id")}} fait toutes les sommes par bâtiment
    sans aucune boucle Python~: pandas regroupe les pixels partageant le même
    identifiant. L'\fonction{id} étant l'index du \fonction{gdf}, le
    \fonction{join} final réaligne tout sans risque.
  \item \textbf{Surface par orientation, pas azimut moyen.} On ne stocke
    \emph{pas} une orientation moyenne -- elle n'aurait aucun sens pour un toit à
    deux pans opposés (la moyenne d'un pan est et d'un pan ouest tomberait au
    nord ou au sud, qu'aucun pan ne regarde). On stocke à la place \textbf{la
    surface inclinée dans chacun des 8 secteurs} (\fonction{surf\_N\_m2},
    \fonction{surf\_E\_m2}, \ldots), bien plus informatif~: « ce bâtiment a tant
    de m$^2$ au sud, tant à l'est ».
  \item \textbf{Pente moyenne sur l'inclinée seulement.} Mélanger les toits plats
    (pente $\approx 0$) fausserait le chiffre~; \fonction{pente\_moy} ne porte
    donc que sur les pixels inclinés.
  \item \textbf{Intensif recalculé à partir de l'extensif.} L'irradiation
    moyenne \fonction{irr\_an\_kwh\_m2} se déduit du total d'énergie divisé par
    la surface totale -- propriété qui rendra la fusion par région triviale
    (section~\ref{sec:afaire})~: on somme les quantités additives (énergies,
    surfaces) puis on redivise.
  \item \textbf{\fonction{join(how="inner")}} ne garde que les bâtiments ayant au
    moins un pixel de toit détecté~; ceux sans toiture exploitable
    (trop bas, mauvaise orientation\ldots) sont écartés du fichier.
\end{itemize}

\subsection{Du gisement à la production : le modèle PV simple}

\fonction{agregerBatiment} ajoute enfin deux colonnes qui passent de l'énergie
\emph{reçue} à l'électricité \emph{produite}, par un modèle volontairement simple
(le niveau « simple » de Buffat)~:

\begin{lstlisting}
res["puissance_kwc"] = res.surf_tot_m2 * TAUX_COUVERTURE * RENDEMENT_MODULE
res["prod_an_kwh"]   = res.irr_an_kwh  * TAUX_COUVERTURE * RENDEMENT_MODULE * PERFORMANCE_RATIO
\end{lstlisting}

\paragraph{Qu'est-ce qu'un kWc ?} Le \textbf{kilowatt-crête} (kWc, ou kWp en
anglais) est la puissance d'une installation PV mesurée dans des
\emph{conditions de test standard} (STC)~: un ensoleillement de $1000$~W/m$^2$,
une température de cellule de $25^\circ$C et un spectre solaire AM1.5, conditions
fixées par la norme IEC~61215\footnote{Définition STC~: voir Réf.~« STC / kWc ».}.
C'est la « taille » nominale d'un système -- le chiffre des fiches techniques --
pas sa production réelle, qui dépend du climat. Un module de rendement $\eta$
soumis à ces $1000$~W/m$^2 = 1$~kW/m$^2$ produit exactement $\eta$~kWc par m$^2$~;
d'où, à couverture pleine, $\text{puissance} = \text{surface}\times\eta$.

\paragraph{Le facteur de production.} La production annuelle est l'irradiation
reçue multipliée par un facteur unique, $\tau\times\eta\times PR \approx 0{,}156$~:
\begin{itemize}
  \item $\eta_{\text{module}}$, \textbf{rendement du panneau} $= 0{,}20$. Valeur
    volontairement \emph{prudente}~: les modules silicium cristallin commerciaux
    atteignaient en moyenne $22{,}0\,\%$ fin 2024 (extrêmes $18{,}8$--$23{,}8\,\%$
    selon le fabricant)\footnote{Fraunhofer ISE, \emph{Photovoltaics Report}, éd.
    2025 -- voir Réf.}, mais un parc de toitures réel comporte des installations
    plus anciennes et moins performantes.
  \item $PR$, \textbf{performance ratio} $= 0{,}78$~: regroupe toutes les autres
    pertes (onduleur, échauffement des cellules, salissures, câblage,
    désadaptation). Valeur dans la fourchette usuelle $0{,}75$--$0{,}85$, le
    « moderne » dépassant $0{,}80$\footnote{\emph{Photovoltaic system
    performance} (synthèse) et IEA-PVPS -- voir Réf.}.
  \item $\tau$, \textbf{taux de couverture} du toit par les panneaux ($1$ par
    défaut~: borne haute, toute la surface exploitable couverte).
\end{itemize}

\paragraph{Puissance et m$^2$/kWc : une conséquence, pas un paramètre.} La
puissance installable découle du \textbf{même} rendement
($\text{puissance} = \text{surface}\times\tau\times\eta$). La « surface par kWc »
classique n'est donc \emph{pas} une constante libre mais s'en déduit~:
$\text{m}^2/\text{kWc} = 1/(\eta\,\tau) \approx 5$ -- cohérent avec la règle
métier de $5$ à $8$~m$^2$/kWc\footnote{Ordre de grandeur installateurs / IRENA --
voir Réf.} (notre $5$ correspond au cas optimiste $\tau=1$). Une seule grandeur de
rendement, deux sorties cohérentes. Les trois constantes ($\eta$, $PR$, $\tau$)
sont en tête de \fichier{irr\_calcul.py}, ajustables d'une ligne. Le modèle
complet à une diode de Buffat (température de cellule, courbe d'onduleur) reste un
raffinement futur (section~\ref{sec:afaire}), qui ne change le résultat que de
quelques pour-cent.

\subsection{\fonction{irradianceTuile} : l'orchestrateur}

\textbf{Rôle.} Enchaîner les trois étapes et écrire le fichier. C'est le point
d'entrée appelé par \fichier{main}.

\begin{lstlisting}
def irradianceTuile(masque_incline, masque_plat, pente, aspect, gdf, horizon, lat, lon, out_path):
    B, D, SAZ, SEL = chargerTable(lat, lon)         # table meteo de la cellule
    df  = irrPixels(masque_incline, masque_plat, pente, aspect,
                    B, D, SAZ, SEL, horizon, ALPHAS, BETAS)
    out = agregerBatiment(df, gdf)
    out.to_file(out_path, driver="GPKG")
    return out
\end{lstlisting}

On charge la table de la cellule (cette fois $SAZ, SEL$ servent à l'ombrage), on
calcule l'irradiance par pixel (ombrage compris), on agrège par bâtiment, on
écrit. L'horizon vient de \fonction{compHZ} (branche B). L'écriture du
GeoPackage est une seule ligne (\fonction{to\_file}) -- pas besoin d'en faire une
fonction à part.

\subsection{Le GeoPackage de sortie}

Une couche vectorielle (driver \fichier{GPKG}), \textbf{une ligne par bâtiment}
ayant au moins un pixel de toiture, en Lambert 93. Colonnes~:

\begin{center}
\begin{tabular}{@{}l p{9cm}@{}}
\toprule
\fichier{cleabs}, \fichier{nature}, \fichier{usage\_1}, & attributs BD TOPO
conservés (clé de jointure + typologie)\\
\fichier{hauteur}, \fichier{nombre\_d\_etages} & \\
\fichier{geometry} & polygone d'emprise du bâtiment\\
\midrule
\fichier{irr\_an\_kwh} & irradiation annuelle \textbf{totale} reçue (kWh/an)\\
\fichier{irr\_an\_kwh\_m2} & irradiation annuelle \textbf{moyenne} (kWh/m$^2$/an)\\
\fichier{prod\_an\_kwh} & production électrique PV estimée (kWh/an)\\
\fichier{puissance\_kwc} & puissance installable estimée (kWc)\\
\fichier{surf\_tot\_m2} & surface de toiture exploitable totale (m$^2$ réels)\\
\fichier{surf\_inclinee\_m2}, \fichier{surf\_plate\_m2} & dont inclinée / plate\\
\fichier{surf\_N\_m2} \ldots \fichier{surf\_NO\_m2} & surface inclinée par
secteur d'orientation (8)\\
\fichier{pente\_moy} & pente moyenne des pans inclinés (degrés)\\
\fichier{nb\_pixels} & nombre de pixels de toit (contrôle qualité)\\
\bottomrule
\end{tabular}
\end{center}

\paragraph{Lire une ligne (exemple réel, dalle DIDIER).} Un bâtiment de
$\fonction{surf\_tot\_m2} = 122$~m$^2$ de toiture exploitable \emph{reçoit}
$\fonction{irr\_an\_kwh} \approx 175\,000$~kWh/an de soleil (soit
$\fonction{irr\_an\_kwh\_m2} \approx 1430$~kWh/m$^2$/an). On pourrait y installer
$\fonction{puissance\_kwc} \approx 24$~kWc de panneaux ($122 \times 0{,}20$), qui
\emph{produiraient} $\fonction{prod\_an\_kwh} \approx 27\,000$~kWh/an
d'électricité ($175\,000 \times 0{,}156$). Les colonnes \fonction{surf\_S\_m2},
\fonction{surf\_E\_m2}\ldots disent vers où sont orientés ces m$^2$ (ici surtout
sud-est). \textbf{À retenir}~: \fonction{irr\_an\_kwh} est l'énergie solaire
\emph{reçue}, \fonction{prod\_an\_kwh} l'électricité \emph{produite}~; ne pas les
confondre.

\paragraph{Validation sur la dalle DIDIER (Poitiers, 0495\_6611).} La pipeline
complète tourne en $\sim 30$~s (dominée par le calcul d'horizon~; la phase tuile
elle-même $< 1$~s). Sortie~: 360 bâtiments. L'irradiation moyenne
\fonction{irr\_an\_kwh\_m2} va de 769 à 1474~kWh/m$^2$/an (moyenne 1319) --
cohérent avec PVGIS pour la région. \textbf{Effet de l'ombrage}~: sans masquage
la moyenne était 1379~; avec l'horizon elle descend à 1319 ($-4{,}4\,\%$) et le
minimum chute de 1214 à 769 (un bâtiment bas très encaissé), ce qui est
physiquement attendu. Contrôles de cohérence vérifiés~: \fonction{surf\_tot} $=$
inclinée $+$ plate, la somme des 8 secteurs égale la surface inclinée, et
\fonction{irr\_an\_kwh\_m2} $=$ \fonction{irr\_an\_kwh} $/$ \fonction{surf\_tot}.
Comme le masque ne garde que les orientations $90$--$270^\circ$, les secteurs N,
NE, NO sont vides (attendu). La logique d'ombrage est aussi validée par un test
unitaire dédié (seuil, limite incluse, choix de la bonne direction d'horizon).
\textbf{Production PV} (modèle simple)~: sur la dalle, $40{,}6$~GWh/an reçus
$\rightarrow$ $6{,}3$~GWh/an produits, pour $6{,}1$~MWc installables.

% ============================================================================
\section{Conventions, unités et pièges récapitulés}
% ============================================================================

\begin{center}
\small
\begin{tabular}{@{}l p{6.4cm} l@{}}
\toprule
\textbf{Grandeur} & \textbf{Unité / convention} & \textbf{Où}\\
\midrule
Coordonnées dalles, rasters & m, Lambert 93 (EPSG 2154) & branche B\\
Coordonnées météo, soleil & degrés, WGS84 (EPSG 4326) & branche A\\
Nom de dalle IGN & km du coin \textbf{nord-ouest} & \fonction{tileBounds}, \fonction{nomCoord}\\
Azimuts (toits, soleil, horizon) & boussole~: 0=N, 90=E, 180=S, 270=O & partout\\
Pente $\beta$ & degrés, 0 = plat & \fonction{clip}, tables\\
Angles d'horizon & \textbf{degrés} (jamais de radians en sortie) & \fonction{compHZ}\\
Heures & \textbf{UTC} (jamais l'heure légale) & branche A\\
Irradiance & W/m$^2$ (tables $B$, $D$) & \fonction{transpAgr}\\
Irradiation & kWh/m$^2$/an & \fonction{totAb}, sortie\\
Surface d'un pixel & $0{,}25/\cos\beta$ m$^2$ (réelle, pas projetée) & \fonction{irrPixels}\\
Secteurs d'orientation & 8 de $45^\circ$, $0=$N, $2=$E, $4=$S, $6=$O & \fonction{agregerBatiment}\\
Nodata rasters & $-9999$, converti en NaN en mémoire & branche B\\
Identifiant bâtiment dans rasters & index GeoDataFrame $+ 1$ (0 = fond) & \fonction{clip}\\
Indices numpy & \fonction{[ligne, colonne]}, ligne 0 au nord & branche B\\
Fichier de sortie & GeoPackage, 1 ligne / bâtiment, Lambert 93 & \fonction{irradianceTuile}\\
\bottomrule
\end{tabular}
\end{center}

Pièges déjà rencontrés et neutralisés, à garder en tête en relisant le code~:
le zéro négatif dans les noms de fichiers (méridien de Greenwich), la
propagation des NaN par \fonction{np.maximum}, \fonction{np.isnan} interdit sur
des entiers, l'ordre (lon, lat) de pyproj, et la division par
$\cos\theta_z$ au crépuscule.

% ============================================================================
\section{Ce qu'il reste à faire}
\label{sec:afaire}
% ============================================================================

La phase tuile (irradiation par pixel avec ombrage, agrégation par bâtiment,
GeoPackage) est \textbf{faite} (section~\ref{sec:tuile}). Restent les
raffinements suivants.

\subsection{Ombrage : ce qui est capté, ce qui ne l'est pas}
\label{sec:ombrage-limites}

L'ombrage \emph{est} appliqué, mais l'horizon est calculé sur la \textbf{dalle
seule}, dans un rayon de $100$~m. Il faut distinguer nettement ce que cela traite
correctement de ce que cela manque -- c'est la principale question de validité de
la version actuelle.

\paragraph{Ce qui est correctement capté.} Pour tout pixel à plus de $100$~m du
bord de la dalle, \emph{tous} les obstacles de son rayon de recherche sont
présents dans la dalle~: son horizon proche est \textbf{complet}. Cela couvre
l'\textbf{auto-ombrage} (un faîtage ou un étage haut qui ombrage une partie basse
du même toit) et l'\textbf{ombrage par les voisins proches} (bâtiments, arbres à
moins de 100~m) -- l'effet d'ombrage dominant en zone bâtie. Le carré central de
$800\times800$~m, soit \textbf{64\,\% de la dalle}, est donc traité correctement.
C'est pourquoi l'ombrage mono-dalle est déjà utile et pas un simple placeholder.

\paragraph{Ce qui est manqué.}
\begin{enumerate}
  \item \textbf{La bande de bord} (l'anneau de 100~m le long des 4 côtés,
    $\sim$36\,\% de la dalle)~: un rayon qui sort de la dalle « retombe » sur le
    pixel de bord (effet du \fonction{np.clip}) au lieu de voir les obstacles de
    la dalle voisine. L'horizon y est donc \emph{sous-estimé}, l'irradiation
    \emph{sur-estimée}. \textbf{Cette perte est irréductible depuis une seule
    dalle}~: l'information n'existe pas dans le fichier, aucune astuce ne peut la
    recréer. La seule vraie correction est de \textbf{charger les dalles voisines}
    (assemblage 3$\times$3, ou blocs glissants pour une région,
    section~\ref{sec:horizon}), à coût quasi nul~: une dalle pèse 16~Mo, le
    3$\times$3 fait 144~Mo ($0{,}016$~Go/km$^2$).
  \item \textbf{Le relief lointain} (collines, montagnes au-delà de $\sim$100~m)~:
    non capté. Buffat le traite par une \emph{seconde} méthode -- l'horizon depuis
    le centre du bâtiment sur un MNT grossier jusqu'à 5~km, puis le maximum des
    deux. Ce terme est ici omis. Il est \textbf{négligeable en plaine} (cas de la
    Vienne) mais compte en vallée encaissée ou en montagne, où un relief peut
    masquer le soleil bas d'hiver plusieurs heures. Donnée mobilisable plus tard~:
    le MNT national IGN \fichier{BD ALTI} à 25~m (ou le \fichier{RGE ALTI}
    rééchantillonné)\footnote{IGN, BD ALTI / RGE ALTI -- voir Réf.}.
  \item \textbf{Le diffus n'est pas masqué} (cf. limites)~: seul le direct est
    éteint à l'ombre.
\end{enumerate}

\paragraph{Que faire sans les voisines ?} On ne peut pas \emph{reconstituer}
l'horizon manquant, mais on peut~: (i) \textbf{marquer} les bâtiments de la bande
de bord ($<100$~m d'un côté) d'un drapeau « horizon incomplet » pour les signaler
dans le cadastre, voire les exclure des statistiques~; (ii) porter le rayon de
$100$ à $500$~m (la valeur de Buffat) pour mieux capter les obstacles
\emph{internes} mi-distants -- au prix de $\times 5$ sur le temps d'horizon, et
\emph{sans} rien changer à la bande de bord. La correction réelle reste le
chargement des voisines.

\subsection{Interpolation bilinéaire $(\alpha, \beta)$}

La phase tuile lit aujourd'hui la case $(\alpha,\beta)$ \textbf{la plus proche}
(plus proche voisin). Une \textbf{interpolation bilinéaire} entre les 4 nœuds
voisins lisserait les sauts d'irradiation entre cases. Deux précautions~: l'axe
des orientations est \textbf{périodique} (interpoler entre $345^\circ$ et
$0^\circ$ via le modulo, sans « traverser » la table), et la pente se borne à
$[0, 70]$. Gain modéré (la grille $15^\circ\times10^\circ$ est déjà fine), à
faire seulement si les cartes paraissent « en escalier ».

\subsection{Profils mensuels et toits plats}

\begin{itemize}
  \item \textbf{Mensuel.} La sortie est aujourd'hui annuelle~; on peut ajouter
    les \textbf{12 valeurs mensuelles} par bâtiment (même intégration sans sommer
    les mois) pour le profil saisonnier.
  \item \textbf{Toits plats.} Les pixels plats reçoivent actuellement
    l'irradiation \emph{horizontale} (leur géométrie réelle, $\beta\approx0$). En
    pratique on y poserait des panneaux \textbf{sur châssis incliné}~; il faudra
    leur attribuer un couple $(\alpha,\beta)$ conventionnel (plein sud à
    l'inclinaison optimale) au moment du modèle PV.
  \item \textbf{Quartiles interannuels.} Comme Buffat, en plus de la moyenne long
    terme, les profils des années de premier et troisième quartiles donnent une
    fourchette~; \fonction{profMH} s'étend facilement (grouper par année d'abord).
\end{itemize}

\subsection{Modèle PV complet (raffinement)}

La conversion en électricité existe déjà en version \textbf{simple}
(section~\ref{sec:tuile})~: production $= E \times \eta_{\text{module}} \times PR$
et puissance $= \text{surface} \times \tau \times \eta$ (donc $\sim 5$~m$^2$ par
kWc, déduit du rendement). C'est suffisant pour un premier cadastre. Reste, si
l'on veut comparer des technologies ou valider sur
des productions réelles, le \textbf{modèle complet} de Buffat (section 3.2)~:
modèle électrique à une diode (paramétrisation de Batzelis), dépendance à la
température de cellule (NOCT), courbe de rendement d'onduleur, dégradation. Il ne
change le résultat que de quelques pour-cent.

\subsection{Passage à l'échelle : multi-dalles et fusion par région}

\begin{itemize}
  \item \textbf{Boucle multi-dalles}~: transformer \fichier{main\_MNTMNS.py} en
    traitement par lot (liste de dalles, téléchargement automatique des MNS/MNT
    via la géoplateforme IGN), avec parallélisation par dalle (chaque dalle est
    indépendante -- \fonction{multiprocessing} trivial). Un \fichier{.gpkg} par
    dalle, comme on fait déjà un \fichier{.npz} par cellule météo.
  \item \textbf{Fusion par région}~: un bâtiment à cheval sur plusieurs dalles a
    ses pixels répartis entre elles, et apparaît donc dans plusieurs
    \fichier{.gpkg} avec des valeurs \emph{partielles}. On ne concatène donc pas~:
    on \fonction{groupby("cleabs")} et on \textbf{somme les colonnes extensives}
    (énergies, surfaces, \fichier{nb\_pixels}), on garde une copie de la géométrie
    et des attributs, puis on \textbf{recalcule les intensives}
    (\fichier{irr\_an\_kwh\_m2} $=$ énergie totale $/$ surface totale). C'est
    exact car énergie et surface sont additives, et les dalles partitionnent
    l'espace (aucun double comptage de pixel).
\end{itemize}

\subsection{Validation}

\begin{itemize}
  \item comparer \fonction{totAb} aux estimations PVGIS en ligne pour plusieurs
    cellules et orientations (fait ponctuellement~: écarts cohérents)~;
  \item comparer les pentes/aspects extraits à des toits connus (contrôle QGIS)~;
  \item à terme, confronter la production simulée à des installations réelles
    (open data Enedis, à la manière des 500 installations de Buffat).
\end{itemize}

\subsection{Limites connues et choix à réévaluer}
\label{sec:limites}

\begin{itemize}
  \item \textbf{Horizon sur la dalle seule}~: ombrage sous-estimé sur la bande de
    bord (36\,\% de la dalle) et relief lointain ignoré -- analyse détaillée
    section~\ref{sec:ombrage-limites}. C'est la limite principale aujourd'hui.
  \item \textbf{Masquage binaire du direct seulement}~: le diffus n'est pas réduit
    par les obstacles (pas de \emph{facteur de vue du ciel})~; or un pixel au fond
    d'une cour « voit » moins de ciel et reçoit donc moins de diffus. Approximation
    de Buffat, conservée pour rester simple~; l'erreur est modérée car en France le
    direct représente la majorité de l'irradiation annuelle.
  \item \textbf{Position solaire moyennée pour l'ombrage}~: le test d'ombre
    utilise l'azimut/élévation \emph{moyens} par (mois, heure) ($SAZ$, $SEL$), pas
    la position instantanée. Un cas où le soleil frôle l'horizon peut être mal
    classé. (La transposition $B$, $D$, elle, utilise les vraies positions
    instantanées avant moyennage~; seul le test d'ombre est moyenné.)
  \item \textbf{Résolution d'horizon $10^\circ$} (36 directions)~: un obstacle fin
    (cheminée, mât) entre deux azimuts échantillonnés peut être manqué. Buffat
    utilise $5^\circ$.
  \item \textbf{Pans nord exclus} par le filtre $90$--$270^\circ$ d'aspect~: assumé
    pour un outil de potentiel « rentable », à lever pour un cadastre exhaustif
    (les secteurs N/NE/NO de sortie sont alors toujours nuls).
  \item \textbf{Toits plats = irradiation horizontale}~: on prend leur géométrie
    réelle ($\beta\approx0$)~; en pratique les panneaux y seraient sur châssis
    inclinés (gain non compté), à traiter au modèle PV.
  \item \textbf{Lecture plus proche voisin} dans la table $(\alpha,\beta)$ (pas
    encore d'interpolation bilinéaire).
  \item \textbf{Albédo constant} ($0{,}2$) et \textbf{pas de neige}~: Buffat note
    un biais hivernal ($+28\,\%$ en décembre) attribué à la neige au sol et sur les
    panneaux~; impact moindre en plaine française.
  \item \textbf{Février = 28 jours} dans \fonction{N\_JOURS}~: biais $< 0{,}1\,\%$
    sur l'année, négligé.
  \item \textbf{Qualité BD TOPO}~: les bâtiments très récents ou non cartographiés
    échappent à la détection~; le MNS, lui, inclut les superstructures (cheminées,
    panneaux existants) qui rugosifient pente et aspect.
\end{itemize}

% ============================================================================
\section*{Références}
\addcontentsline{toc}{section}{Références}
% ============================================================================

\begin{itemize}
  \item R.~Buffat, S.~Grassi, M.~Raubal, \emph{A scalable method for estimating
    rooftop solar irradiation potential over large regions}, Applied Energy 216
    (2018) 389--401. -- La méthode de référence de toute la pipeline.
  \item R.~Perez et al., \emph{Modeling daylight availability and irradiance
    components from direct and global irradiance}, Solar Energy 44 (1990). --
    Le modèle de transposition du diffus.
  \item I.~Reda, A.~Andreas, \emph{Solar position algorithm for solar radiation
    applications}, Solar Energy 76 (2004). -- L'algorithme SPA.
  \item PVGIS~5.3, Commission européenne / JRC~:
    \url{https://re.jrc.ec.europa.eu/pvg_tools/fr/} (API \fichier{.../api/v5\_3/}).
    -- Données SARAH-3.
  \item pvlib python~: \url{https://pvlib-python.readthedocs.io} --
    bibliothèque de référence pour le solaire (positions, transposition,
    iotools PVGIS).
  \item IGN, LiDAR HD, BD TOPO, BD ALTI et RGE ALTI~:
    \url{https://geoservices.ign.fr}. -- Données d'entrée et MNT national
    (BD ALTI 25~m) mobilisable pour l'horizon lointain.
  \item \textbf{SARAH-3} (CM SAF / EUMETSAT)~: U.~Pfeifroth et al.,
    \emph{SARAH-3 -- satellite-based climate data records of surface solar
    radiation}, Earth Syst. Sci. Data 16 (2024). -- Rayonnement satellite
    ($0{,}05^\circ$, pas 30~min, depuis 1983) derrière PVGIS.
  \item \textbf{STC / kWc}~: norme IEC~61215~; synthèse
    \url{https://en.wikipedia.org/wiki/Nominal_power_(photovoltaic)}. -- Définition
    de la puissance-crête.
  \item \textbf{Performance ratio}~: \emph{Photovoltaic system performance}
    (Wikipedia) et IEA-PVPS Task~2. -- Fourchette $0{,}75$--$0{,}85$.
  \item \textbf{Rendement modules}~: Fraunhofer ISE, \emph{Photovoltaics Report}
    (éd. 2025), \url{https://www.ise.fraunhofer.de}. -- Moyenne c-Si
    $22{,}0\,\%$ fin 2024.
  \item \textbf{Surface par kWc}~: ICLEI / IRENA, \emph{Solar PV Rooftop Systems}~;
    règle métier $5$--$8$~m$^2$/kWc.
\end{itemize}

\end{document}
